\section{全模群情形的维数公式}
\label{sec:dim-sl2z}
% ============================================================================

如果只看 $\Gamma=\SL_2(\Z)$，维数公式美得几乎像谜语：大多数偶权空间的维数都只比 $\lfloor k/12\rfloor$ 多一个单位，但偏偏在 $k\equiv 2\pmod{12}$ 时又少掉这一维。我们现在要做的，就是把这条规律完整证明出来。

核心思路其实只有两步。第一步，用 valence formula 控制一个模形式总共有多少零点；这会立刻告诉我们“小权时根本容不下 cusp form”。第二步，用权 $12$ 的 cusp form $\Delta$ 把高权尖点形式降到低 $12$ 个权的模形式，从而得到一个递推。

% figures/ch03-dimMk-staircase.tex
% Staircase for dim M_k(SL_2(Z))
\begin{figure}[ht]
\centering
\begin{tikzpicture}[x=0.55cm,y=0.8cm,font=\small]
  \draw[->] (0,0) -- (31,0) node[right] {$k$};
  \draw[->] (0,0) -- (0,4.3) node[above] {$\dim \Mk_k(\SL_2(\Z))$};

  \foreach \x in {0,4,6,8,10,12,14,16,18,20,22,24,26,28,30}
    \draw (\x,0) -- (\x,-0.08) node[below] {\x};
  \foreach \y in {1,2,3,4}
    {\draw[gray!25] (0,\y) -- (30.5,\y); \node[left] at (0,\y) {\y};}

  \draw[very thick,blue!70!black]
    (0,1) -- (2,1) -- (4,1) -- (6,1) -- (8,1) -- (10,1)
    -- (12,2) -- (14,1) -- (16,2) -- (18,2) -- (20,2) -- (22,2)
    -- (24,3) -- (26,2) -- (28,3) -- (30,3);

  \foreach \x/\y in {0/1,4/1,6/1,8/1,10/1,12/2,14/1,16/2,18/2,20/2,22/2,24/3,26/2,28/3,30/3}
    \fill[blue!70!black] (\x,\y) circle (2pt);

  \node[align=left,fill=white] at (22,3.7) {遇到 $k\equiv 2\pmod{12}$\\会“少一格”};
\end{tikzpicture}
\caption{$\dim \Mk_k(\SL_2(\Z))$ 的“楼梯图”。大体规律是每增加 $12$ 个权，维数增加 $1$；但在 $k\equiv 2 \pmod{12}$ 时，维数比朴素的“$\lfloor k/12\rfloor+1$”少一。}
\label{fig:dimMk-staircase}
\end{figure}


\begin{theorem}[全模群的维数公式]
\label{thm:dim-sl2z}
对 $\SL_2(\Z)$，有：
\begin{enumerate}
  \item 若 $k<0$ 或 $k$ 为奇数，则
        \[
          \dim \Mk_k(\SL_2(\Z))=0.
        \]
  \item $\dim \Mk_0(\SL_2(\Z))=1$，而 $\dim \Mk_2(\SL_2(\Z))=0$。
  \item 若 $k\ge 4$ 为偶数，则
        \[
          \dim \Mk_k(\SL_2(\Z))=
          \begin{cases}
            \left\lfloor\dfrac{k}{12}\right\rfloor+1,& k\not\equiv 2 \pmod{12},\\[6pt]
            \left\lfloor\dfrac{k}{12}\right\rfloor,& k\equiv 2 \pmod{12}.
          \end{cases}
        \]
  \item 若 $k\ge 4$ 为偶数，则
        \[
          \dim \Sk_k(\SL_2(\Z))=\dim \Mk_k(\SL_2(\Z))-1,
        \]
        而 $\dim \Sk_2(\SL_2(\Z))=0$。
\end{enumerate}
\end{theorem}

\begin{example}[先看两个最重要的数]
\label{ex:dim-m12-s24}
由定理立刻得到
\[
  \dim \Mk_{12}(\SL_2(\Z))=2,
  \qquad
  \dim \Sk_{24}(\SL_2(\Z))=2.
\]
前者意味着权 $12$ 模形式由一条 Eisenstein 方向和一条 cusp direction 组成；事实上
\[
  \Mk_{12}(\SL_2(\Z))=\C E_{12}\oplus \C\Delta.
\]
后者则说明权 $24$ 的尖点形式已经不是“自动特征形式”的一维世界了；例如
\[
  \Sk_{24}(\SL_2(\Z))=\Delta\Mk_{12}(\SL_2(\Z))
  =\C(\Delta E_{12})\oplus \C\Delta^2.
\]
真正想找到 Hecke 特征形式时，就必须在这个二维空间里对角化算子。
\end{example}

\begin{lemma}[valence formula]
\label{lem:valence-formula}
设 $f\in \Mk_k(\SL_2(\Z))$ 为非零模形式，权 $k\ge 0$。则
\[
  \ord_\infty(f)+\frac{1}{2}\ord_i(f)+\frac{1}{3}\ord_\rho(f)
  +\sum_{\tau\in \SL_2(\Z)\backslash\HH\setminus\{i,\rho\}} \ord_\tau(f)
  =\frac{k}{12}.
\]
这里 $\rho=e^{2\pi i/3}$，而求和按轨道各取一个代表。
\end{lemma}

\begin{example}[用 valence formula 立刻看出 $\Delta$ 的特殊性]
\label{ex:valence-delta}
$\Delta$ 是权 $12$ 的 cusp form，在 $\infty$ 处恰有一阶零点，而且它在 $\HH$ 上无零点。于是 valence formula 左边只有一项贡献：
\[
  \ord_\infty(\Delta)=1=\frac{12}{12}.
\]
这正说明 $\Delta$ 是“最小可能”的非零 cusp form，也解释了为什么它在维数递推里会扮演基本单位的角色。
\end{example}

\begin{proof}
\textbf{第 1 步：奇权必为零。}
若 $k$ 为奇数，而 $f\in \Mk_k(\SL_2(\Z))$，则因为 $-I\in\SL_2(\Z)$，有
\[
  f(\tau)=f((-I)\tau)=(-1)^k f(\tau)=-f(\tau).
\]
故 $f=0$。同理，负权全纯模形式也不可能存在，因此 $k<0$ 时维数为零。

\textbf{第 2 步：把 $\Mk_k$ 分成 Eisenstein 部分与 cusp 部分。}
对每个偶数 $k\ge 4$，第二章已构造出非零 Eisenstein 级数 $E_k\in \Mk_k(\SL_2(\Z))$，其常数项为 $1$，因此不是 cusp form。于是
\[
  \Mk_k(\SL_2(\Z))=\Sk_k(\SL_2(\Z))\oplus \C E_k,
\]
从而
\[
  \dim \Mk_k=\dim \Sk_k+1 \qquad (k\ge 4\text{ even}).
\]

\textbf{第 3 步：小于 $12$ 时不存在非零 cusp form。}
设 $0<k<12$ 为偶数，且 $f\in \Sk_k(\SL_2(\Z))$ 非零。因为 $f$ 是 cusp form，
\[
  \ord_\infty(f)\ge 1.
\]
另一方面，valence formula 给出
\[
  \ord_\infty(f)\le \frac{k}{12}<1,
\]
矛盾。因此
\[
  \Sk_k(\SL_2(\Z))=0 \qquad (0<k<12,\ k\text{ even}).
\]
于是
\[
  \dim \Mk_k(\SL_2(\Z))=1 \quad (k=4,6,8,10),
\]
而 $k=2$ 时既没有非零 cusp form，也没有真正的 Eisenstein 级数，所以
\[
  \dim \Mk_2(\SL_2(\Z))=0.
\]

\textbf{第 4 步：乘上 $\Delta$ 给出递推。}
对任意偶数 $k\ge 12$，考虑映射
\[
  \Phi\colon \Mk_{k-12}(\SL_2(\Z))\longrightarrow \Sk_k(\SL_2(\Z)),
  \qquad g\longmapsto \Delta g.
\]
它显然线性，且因为 $\Delta$ 是权 $12$ cusp form，故像落在 $\Sk_k$ 中。

这个映射还是双射。单射显然，因为 $\Delta\neq 0$。为证满射，取任意 $f\in \Sk_k(\SL_2(\Z))$。由 \cref{prop:E4-E6-Delta}，$\Delta$ 在 $\HH$ 上无零点，并且在 $\infty$ 处恰有一阶零点；又因为 $f$ 在 $\infty$ 至少有一阶零点，所以
\[
  \frac{f}{\Delta}
\]
在 $\HH$ 和尖点处都全纯。其变换律权重为 $k-12$，故 $f/\Delta\in \Mk_{k-12}(\SL_2(\Z))$。于是
\[
  \Sk_k(\SL_2(\Z))\cong \Mk_{k-12}(\SL_2(\Z))
  \qquad (k\ge 12\text{ even}).
\]
因此
\[
  \dim \Mk_k=\dim \Sk_k+1=\dim \Mk_{k-12}+1
  \qquad (k\ge 12\text{ even}).
\]

\textbf{第 5 步：解递推。}
写 $k=12m+r$，其中 $r\in\{0,2,4,6,8,10\}$。递推反复使用 $m$ 次后，得到
\[
  \dim \Mk_k=\dim \Mk_r+m.
\]
而基例是
\[
  \dim \Mk_0=1,\quad \dim \Mk_2=0,\quad
  \dim \Mk_4=\dim \Mk_6=\dim \Mk_8=\dim \Mk_{10}=1.
\]
于是：若 $r=2$，就有 $\dim \Mk_k=m=\lfloor k/12\rfloor$；其余偶数余类都给出
\[
  \dim \Mk_k=m+1=\left\lfloor\frac{k}{12}\right\rfloor+1.
\]
最后再用
\[
  \dim \Sk_k=\dim \Mk_k-1 \qquad (k\ge 4\text{ even})
\]
即可得到尖点形式空间的公式。
\end{proof}

\begin{remark}[与 $E_4,E_6$ 的关系]
从维数公式还能反过来读出一个非常重要的结论：在每个偶权下，$\Mk_k(\SL_2(\Z))$ 都由满足
\[
  4a+6b=k
\]
的单项式 $E_4^aE_6^b$ 张成，而且这些单项式的个数正好等于上面的维数。这将在习题里再次出现：对 $k<12$，每个模形式都只能是 $E_4,E_6$ 的一个多项式。
\end{remark}

\begin{theorem}[全模群模形式的分次环]
\label{thm:graded-ring-sl2z}
把所有偶权模形式放在一起，得到分次环
\[
  \Mk_*(\SL_2(\Z))=\bigoplus_{k\ge 0}\Mk_k(\SL_2(\Z)).
\]
则自然映射
\[
  \C[X,Y]\longrightarrow \Mk_*(\SL_2(\Z)),
  \qquad X\longmapsto E_4,
  \quad Y\longmapsto E_6
\]
给出分次环同构。
换言之，
\[
  \Mk_*(\SL_2(\Z))=\C[E_4,E_6],
\]
并且 $E_4,E_6$ 代数无关。
\end{theorem}

\begin{proof}
记上面的分次环同态为 $\phi$。
先证它满射。
对权 $k$ 作归纳。
当 $k=0,4,6,8,10$ 时，维数公式表明 $\Mk_k(\SL_2(\Z))$ 都是一维，
相应生成元可取为 $1,E_4,E_6,E_4^2,E_4E_6$；
而 $\Mk_2(\SL_2(\Z))=0$。

现设 $k\ge 12$ 且结论对更低偶权成立，取任意
$f\in \Mk_k(\SL_2(\Z))$。
因为总能找到某个满足 $4a+6b=k$ 的单项式 $P_k(E_4,E_6)$，
且它的常数项为 $1$，
故
\[
  f-a_0(f)P_k(E_4,E_6)
\]
在 $\infty$ 处常数项为零，所以是一个 cusp form。
由前面的递推，尖点形式都可唯一写成 $\Delta g$，其中
$g\in \Mk_{k-12}(\SL_2(\Z))$。
按归纳假设，$g$ 已是 $E_4,E_6$ 的多项式，
而 $\Delta=(E_4^3-E_6^2)/1728$ 也是 $E_4,E_6$ 的多项式，
故 $f$ 也是如此。
于是 $\phi$ 满射。

再证 $E_4,E_6$ 代数无关。
固定权 $k$，考察所有满足 $4a+6b=k$ 的单项式 $E_4^aE_6^b$。
这样的 $(a,b)$ 个数恰好等于 \cref{thm:dim-sl2z} 给出的 $\dim \Mk_k(\SL_2(\Z))$。
而上一步已经说明这些单项式张成整个 $\Mk_k(\SL_2(\Z))$，
故它们实际上构成一组基。
因此同一权里不存在非平凡线性关系，
更不可能存在非零多项式关系 $P(E_4,E_6)=0$。
于是 $\phi$ 既满又无核，从而是同构。
\end{proof}

\begin{corollary}[尖点形式理想由 $\Delta$ 生成]
\label{cor:cusp-ideal-generated-by-delta}
在分次环 $\Mk_*(\SL_2(\Z))$ 中，尖点形式组成的分次理想满足
\[
  \Sk_*(\SL_2(\Z))=\Delta\cdot \Mk_{*-12}(\SL_2(\Z)).
\]
也就是说，每个 cusp form 都唯一地写成 $\Delta g$，其中 $g$ 是较低权的模形式。
\end{corollary}

\begin{proof}
若 $g\in \Mk_{k-12}(\SL_2(\Z))$，则 $\Delta g$ 显然是权 $k$ 的 cusp form。
反过来，若 $f\in \Sk_k(\SL_2(\Z))$，则 $f$ 在 $\infty$ 至少有一阶零点，
而 $\Delta$ 在 $\infty$ 恰有一阶零点、在 $\HH$ 上无零点，
故 $f/\Delta$ 在整条模曲线上全纯，属于 $\Mk_{k-12}(\SL_2(\Z))$。
唯一性来自 $\Delta\neq 0$。
\end{proof}

\begin{corollary}[维数差一的简洁证明]
\label{cor:dimMk-dimSk-plus-one-short}
对偶数权 $k\ge 4$，有
\[
  \dim \Mk_k(\SL_2(\Z))=\dim \Sk_k(\SL_2(\Z))+1.
\]
而在 $k=2$ 时，两边都等于 $0$。
\end{corollary}

\begin{proof}
取任意满足 $4a+6b=k$ 的单项式 $P_k(E_4,E_6)$。
因为它的常数项为 $1$，所以它不在 $\Sk_k$ 中。
另一方面，\cref{cor:cusp-ideal-generated-by-delta} 表明每个模形式都可唯一写成
\[
  f=c\,P_k(E_4,E_6)+\Delta g,
  \qquad c\in \C,
  \quad g\in \Mk_{k-12}(\SL_2(\Z)).
\]
于是
\[
  \Mk_k(\SL_2(\Z))=\C\,P_k(E_4,E_6)\oplus \Sk_k(\SL_2(\Z)),
\]
立得维数差一。
$k=2$ 时没有非零模形式，所以是唯一例外。
\end{proof}
